\documentclass[11pt]{article}
\usepackage{shared/template_tgir_article}
\usepackage{booktabs}
\usepackage{longtable}
\usepackage{tabularx}
\usepackage{array}
\usepackage{enumitem}
\usepackage{xcolor}
\usepackage{listings}

\lstset{
    basicstyle=\ttfamily\small,
    breaklines=true,
    columns=fullflexible,
    frame=single
}

\newcommand{\product}{TGIR QuantLib Tools}
\newcommand{\code}[1]{\texttt{#1}}
\newcommand{\docdate}{2026-04-06}

% Visual design is controlled exclusively by template_tgir_article.sty and
% IEEEtran. Keep this file limited to content-support packages and commands.


\title{Quant Guide\\\large \product}
\author{TG Investments and Research}
\date{\docdate}

\begin{document}
\maketitle
\tableofcontents
\newpage

\section{Scope}
This guide is for quantitative users who need to understand the pricing assumptions, calibration logic, risk measures, and research references used in the workstation.

\section{Portfolio Scope}
The repo prices five instruments:
\begin{itemize}[leftmargin=1.5em]
    \item a spot-starting fixed-versus-SOFR swap,
    \item a European swaption under a normal-volatility convention priced with Hull-White 1F plus Jamshidian's decomposition,
    \item a Bermudan swaption under a calibrated short-rate model,
    \item a second Bermudan swaption matching the workbook benchmark trade, and
    \item an SPX cliquet under a Black-Scholes-Merton process with analytic cliquet valuation.
\end{itemize}

\section{Rates Curve Construction}
\subsection{Inputs}
The rates market state follows the workbook strip and a configurable valuation date, defaulting to \code{2026-03-10}:
\begin{center}
\begin{tabular}{lllllllllll}
\toprule
\code{1D} & \code{1W} & \code{2W} & \code{3W} & \code{1M} & \code{2M} & \code{3M} & \code{6M} & \code{9M} & \code{1Y} & \code{\dots 30Y} \\
\bottomrule
\end{tabular}
\end{center}

\subsection{Bootstrap}
The curve is constructed in \path{portfolio.py} with:
\begin{itemize}[leftmargin=1.5em]
    \item one \code{ql.DepositRateHelper} for the \code{1D} point, and
    \item a strip of \code{ql.OISRateHelper} objects for the remaining term SOFR OIS points.
\end{itemize}
The term structure is a \code{ql.PiecewiseLogCubicDiscount}. Calibration quality is checked by rebuilding quoted OIS swaps and asserting near-zero NPV.

\section{European Swaption Model}
The European swaption calibrates \code{ql.HullWhite} to the selected ATM normal-volatility matrix pillar and prices with \code{ql.JamshidianSwaptionEngine}. The matrix follows the workbook surface, with short expiries from \code{1M} onward, annual expiries through \code{10Y}, then \code{12Y}, \code{15Y}, \code{20Y}, and \code{25Y}, against swap tenors through \code{30Y}. The on-screen matrix uses the exact workbook axes. When the Bermudan calibration needs a \code{3Y} expiry, it is linearly interpolated between the workbook \code{2Y} and \code{4Y} rows.

\section{Bermudan Swaption Model}
\subsection{Calibration}
The Bermudan stack:
\begin{enumerate}[leftmargin=1.5em]
    \item identifies the trade's Bermudan call schedule and the remaining swap tenor implied by the common fixed final maturity,
    \item creates one \code{ql.SwaptionHelper} per diagonal pillar,
    \item calibrates either \code{ql.HullWhite} or \code{ql.G2} while holding the mean-reversion input fixed, and
    \item prices the Bermudan with \code{ql.TreeSwaptionEngine}.
\end{enumerate}
All rates trades are modeled as \code{ql.OvernightIndexedSwap} instruments with daily compounded SOFR coupons over the chosen reset period. Payment frequency and reset frequency are explicit trade inputs; the workbook reference trade uses annual pay and annual reset.
For example, a \code{5Y NC 2Y} Bermudan keeps the same final maturity date and therefore exercises into the remaining \code{3Y}, \code{2Y}, and \code{1Y} swaps along its annual call schedule.
The workstation defaults Bermudans to Hull-White 1F, and the detail page allows a switch to G2++ so users can compare values and risk outputs under the same market data.

\subsection{Interpretation}
The calibration helpers are labeled by the booked Bermudan schedule, for example \code{1Y x 4Y}, \code{2Y x 3Y}, \code{3Y x 2Y}, \code{4Y x 1Y} for a five-year annual Bermudan. When a booked expiry is not on the editable matrix axes, the helper linearly interpolates in the expiry dimension and keeps the exact source matrix points for audit purposes.

\section{SPX Cliquet Model}
\subsection{State Variables}
The cliquet reads:
\begin{itemize}[leftmargin=1.5em]
    \item SPX spot,
    \item flat dividend yield,
    \item flat Black volatility,
    \item quantity,
    \item percentage moneyness,
    \item maturity in months, and
    \item reset frequency in months.
\end{itemize}

\subsection{Process And Engine}
The market process is a \code{ql.BlackScholesMertonProcess} with:
\begin{itemize}[leftmargin=1.5em]
    \item the shared SOFR curve used as the risk-free curve,
    \item a flat dividend curve, and
    \item a flat Black volatility surface.
\end{itemize}
The instrument is a \code{ql.CliquetOption} with \code{ql.PercentageStrikePayoff} and \code{ql.AnalyticCliquetEngine}.

\subsection{Extra Analytics}
The cliquet page supplements the QuantLib price with:
\begin{itemize}[leftmargin=1.5em]
    \item reset-period decomposition into forward-start option values,
    \item deterministic spot-volatility scenarios,
    \item trade-level analytic Greeks, and
    \item Monte Carlo payoff-distribution diagnostics.
\end{itemize}

\section{Risk Outputs}
\subsection{Rates Trades}
The swap and swaptions expose on-demand point sensitivities:
\begin{itemize}[leftmargin=1.5em]
    \item curve ladder from single-point SOFR quote bumps,
    \item European swaption vega from a single matrix-cell bump,
    \item Bermudan matrix sensitivity from the trade schedule's helper source points, and
    \item a separate short-rate mean-reversion sensitivity for European and Bermudan swaptions.
\end{itemize}

\subsection{Cliquet Trade}
The cliquet page reports:
\begin{itemize}[leftmargin=1.5em]
    \item delta, gamma, vega, theta, rho, and dividend rho,
    \item scenario NPVs for spot and volatility shocks, and
    \item distribution metrics such as 5th percentile PV and expected shortfall.
\end{itemize}

\section{Validation}
The test suite validates the pricing stack in three distinct ways:
\begin{itemize}[leftmargin=1.5em]
    \item calibration repricing for the SOFR OIS curve,
    \item calibration repricing for Bermudan helpers implied by the fixed-maturity call schedule, and
    \item ten cliquet identity cases where the payoff reduces to simpler known structures.
\end{itemize}

\section{Research References}
The curated references live in \path{docs/RESEARCH.md} and on the in-app model page. The swaption list emphasizes:
\begin{itemize}[leftmargin=1.5em]
    \item Peter Caspers' QuantLib-linked swaption implementation notes,
    \item Jamshidian's bond-option decomposition, and
    \item classic Bermudan exercise papers.
\end{itemize}
The cliquet list covers numerical methods, hedging, Levy models, regime-switching, and counterparty-risk extensions.

\section{Model Boundaries}
This workstation deliberately avoids:
\begin{itemize}[leftmargin=1.5em]
    \item live data feeds,
    \item smile or skew modeling for the cliquet,
    \item short-rate models beyond the implemented Hull-White 1F and G2++ calibration paths,
    \item database-backed trade lifecycle, and
    \item portfolio aggregation beyond the five-trade sandbox.
\end{itemize}

\end{document}
